\documentclass[11pt,a4paper]{article}
\usepackage[T1]{fontenc}
\usepackage[utf8]{inputenc}
\usepackage{lmodern}
\usepackage{amsmath,amssymb,amsthm,mathtools}
\usepackage{geometry}
\geometry{margin=2.6cm}
\usepackage{booktabs,array,tabularx}
\usepackage{enumitem}
\usepackage[dvipsnames]{xcolor}
\usepackage[hidelinks]{hyperref}
\usepackage{xurl}
\hypersetup{pdftitle={Bridge and size: the Kodaira-Spencer morphism, arithmetic derivatives and Step (xi) of IUT as lax
2-cells},pdfauthor={Toshihisa Urahigashi}}
\setlist{itemsep=2pt,topsep=4pt}
\raggedbottom
\emergencystretch=3em
\numberwithin{table}{section}

\newcommand{\Z}{\mathbb{Z}}
\newcommand{\Q}{\mathbb{Q}}
\newcommand{\Vbad}{\mathbb{V}^{\mathrm{bad}}}
\newcommand{\qq}{\underline{\underline{q}}}
\newcommand{\Ocal}{\mathcal{O}}
\newcommand{\Fcal}{\mathcal{F}}
\newcommand{\Lcal}{\mathcal{L}}
\newcommand{\Ecal}{\mathcal{E}}
\newcommand{\Tcal}{\mathcal{T}}
\newcommand{\FU}{\mathcal{F}^{\vdash\times\boldsymbol{\mu}}}
\newcommand{\FV}{\mathcal{F}^{\Vdash\blacktriangleright\times\boldsymbol{\mu}}}
\newcommand{\ThetaLGP}{\Theta^{\times\boldsymbol{\mu}}_{\mathrm{LGP}}}
\newcommand{\KS}{\mathrm{KS}}
\newcommand{\leanus}{\textunderscore\allowbreak}
\newcommand{\lean}[1]{{\ttfamily\let\_\leanus #1}}

% labels
\newcommand{\lab}[2]{{\upshape\fcolorbox{#1}{#1!8}{\scriptsize\sffamily\bfseries\textcolor{#1}{#2}}}}
\newcommand{\tagS}{\lab{blue}{S}}
\newcommand{\tagX}{\lab{Purple}{X}}
\newcommand{\tagP}{\lab{ForestGreen}{P}}
\newcommand{\tagR}{\lab{TealBlue}{R}}
\newcommand{\tagH}{\lab{BurntOrange}{H}}
\newcommand{\q}[1]{``#1''}
\newcommand{\srcquote}[2]{\begin{quote}\small #1\par\nopagebreak\hfill{\footnotesize #2}\end{quote}}
\newcommand{\known}[1]{\par{\footnotesize\itshape #1}\par}

\theoremstyle{plain}
\newtheorem{proposition}{Proposition}[section]
\theoremstyle{definition}
\newtheorem{definition}[proposition]{Definition}
\newtheorem{observation}[proposition]{Observation}

% written by make_generated.py from ../scripts/bridge_lib.py; do not edit by hand
\newcommand{\thedoi}{10.5281/zenodo.23984201}
\newcommand{\tabletriples}{%
$2+3^{10}\cdot 109=23^{5}$ & 1.6299 & $3^{10}$ & $b$ & $a/c$ & 10 & 2, 23 \\
 &  & $23^{5}$ & $c$ & $-a/b$ & 5 & 2, 3, 109 \\
$11^{2}+3^{2}\cdot 5^{6}\cdot 7^{3}=2^{21}\cdot 23$ & 1.6260 & $11^{2}$ & $a$ & $b/c$ & 2 & 2, 3, 5, 7, 23 \\
 &  & $3^{2}$ & $b$ & $a/c$ & 2 & 2, 11, 23 \\
 &  & $5^{6}$ & $b$ & $a/c$ & 6 & 2, 11, 23 \\
 &  & $7^{3}$ & $b$ & $a/c$ & 3 & 2, 11, 23 \\
$1+2\cdot 3^{7}=5^{4}\cdot 7$ & 1.5679 & $3^{7}$ & $b$ & $a/c$ & 7 & 5, 7 \\
 &  & $5^{4}$ & $c$ & $-a/b$ & 4 & 2, 3 \\
$283+5^{11}\cdot 13^{2}=2^{8}\cdot 3^{8}\cdot 17^{3}$ & 1.5808 & $5^{11}$ & $b$ & $a/c$ & 11 & 2, 3, 17, 283 \\
 &  & $13^{2}$ & $b$ & $a/c$ & 2 & 2, 3, 17, 283 \\
 &  & $3^{8}$ & $c$ & $-a/b$ & 8 & 5, 13, 283 \\
 &  & $17^{3}$ & $c$ & $-a/b$ & 3 & 5, 13, 283 \\}
\newcommand{\tabledata}{%
D211 & 7 & $211+$ & 58 & 29 & 87 & 232 ($j=3$) & 559.56 & 79.94 \\
D109 & 7 & $109+$ & 58 & 29 & 87 & 232 ($j=3$) & 490.60 & 70.09 \\
D2 & 7 & a 50-digit prime ($-$) & 2 & 1 & 3 & 8 ($j=3$) & 440.17 & 62.88 \\
D3 & 7 & $30965771-$ & 60 & 30 & 90 & 240 ($j=3$) & 2,068.45 & 295.49 \\
D18 & 7 & $30965771+$ & 36 & 18 & 54 & 144 ($j=3$) & 1,127.41 & 161.06 \\
D15 & 7 & $30965819+$ & 30 & 15 & 45 & 120 ($j=3$) & 946.12 & 135.16 \\
D2357$^{\dagger}$ & 157 & $2357+$ & 944 & 472 & 1,416 & 2,871,176 ($j=78$) & 3,840,837.03 & 24,463.93 \\}
\newcommand{\dQ}{21.8011}
\newcommand{\dwall}{79.9372}
\newcommand{\dgap}{559.5605}
\newcommand{\dgaptwo}{87}
\newcommand{\dgapthree}{232}
\newcommand{\ddepth}{29}
\newcommand{\dordq}{58}
\newcommand{\dnbad}{5}


\title{Bridge and size: the Kodaira--Spencer morphism,\\ arithmetic derivatives and Step (xi) of IUT\\ as lax 2-cells}
\author{Toshihisa Urahigashi\\\small ORCID 0009-0004-0460-6242\\\small\texttt{123026.urahigashi.toshihisa@gmail.com}}
\date{Version 1.1.0, 8 November 2026\\\small\href{https://doi.org/\thedoi}{doi:\thedoi}}

\begin{document}
\maketitle

\begin{abstract}
In function fields, the height inequality (the analogue of Szpiro's inequality) follows from the Kodaira--Spencer
morphism. This note places this classical argument, the ``bridges'' that are built in number fields from an additive
relation $a+b=c$ (the $p$-adic logarithms of Baker's method and the arithmetic derivatives of Pasten), and Step~(xi)
of the proof of Corollary~3.12 of inter-universal Teichm\"uller theory (IUT)~III side by side in the form of the lax
2-cells of the earlier note \cite{SL}. Mochizuki himself describes the Kodaira--Spencer morphism as analogous to the
``license to confuse'' of IUT \cite{Alien}. This note is an essay that proposes a viewpoint; it contains almost no new
mathematical results.

There are two points. (1) The difficulty splits into a \emph{bridge} (that a 2-cell exists) and a \emph{size} (the
blurring that must be allowed). This restates, in the language of \cite{SL}, two requirements stated by Mochizuki
himself: a common container that is large enough to house both sides, and small enough to yield a nontrivial estimate
\cite{Roots}. In function fields both come automatically from the structure. A bridge from an additive relation exists
unconditionally, but bounding its size is equivalent to the abc conjecture (Pasten \cite{Pa}). In IUT, conversely, the
blurring is small by the upper bound of IUT~IV, while the bridge (the inclusion in Step~(xi), called a bridge earlier
by Higuchi, Pajo and Geiger \cite{Hi,Pj,Ge3}) is not derived in the written argument. (2) A bridge from an additive relation does not cross the $\Theta$-link. The $\Theta$-link carries the
unit portions by full poly-isomorphisms, so that specific global elements (such as $1-\lambda$ of the Legendre form)
are not carried. If one forces them across, a contradiction appears at the level of units already on actual initial
$\Theta$-data (at the place of the datum D211 above $211$, discrepancies of $\dgaptwo$ and $\dgapthree$ levels for the
labels $2$ and $3$). To absorb this discrepancy by blurring requires exactly half of the discrepancy at the level of
value groups, which in the normalization of Corollary~3.12 (the $q$ side is measured by $\qq=q^{1/2\ell}$) is $\ell$
times the wall $(k-1)Q$. This grows with $q$, while the upper bound $C$ of IUT~IV does not: on the six completely
factored data it is still below $C$ (between $0.36$ and $0.92$ of $C$), and after replacing $q$ by $q^N$ it exceeds $C$
for $N\ge3$ (checked for $N\le40$ \cite{TC}). Moreover, the curves on
which isogenies give a bridge (curves with a global multiplicative subspace) are excluded from the initial
$\Theta$-data by the condition of Definition~3.1~(c) of IUT~I. In summary, the reading of this note is that the ring
structure (in particular, the addition) that IUT gives up in order to circumvent the absence of a global multiplicative
subspace is precisely the material of a Kodaira--Spencer-type bridge (\S8.1).

Hence, to complete IUT as an argument of Kodaira--Spencer type, one would have to build, inside the multiradial
representation, a bridge whose blurring does not grow with $q$. For a single datum this claim is equivalent to the
height inequality \cite{SL}, and if it holds for all data the abc conjecture follows \cite{lana} (for arithmetic
derivatives, the existence of small ones is equivalent to the abc conjecture \cite{Pa}). The language of homotopy type
theory (HoTT) can state this precisely, but cannot construct it. What HoTT contributed here is that, by requiring one to
say along which path one identifies, it made visible that an additive bridge is information on the side of structure
(\S7.4). \textbf{No claim is made about whether IUT is correct.}
\end{abstract}

\tableofcontents

\section{Introduction}

\subsection{The question}
\tagP{} In earlier notes we saw that the dispute over Corollary~3.12 of IUT~III concentrates on the point, in
Step~(xi) of its proof, that the value on the $q$ side is one of the possible values on the $\Theta$ side (below, ``the
inclusion''), and that this point is of the same strength as the height inequality. In the language of \cite{SL},
Corollary~3.12 (in the model) is equivalent to a lax 2-cell between two paths, which is equivalent to $(k-1)Q\le E$
\cite[Theorem~5.3]{SL}. Here $Q=|\log\qq|$, $k=\ell(\ell+1)/12$ is the mean of $j^2$ over the labels
$j=1,\dots,\ell^*$ (where $\ell^*=(\ell-1)/2$), and $E$ is the blurring allowed in the comparison. We call $(k-1)Q$
the \emph{wall}, after the word that the interim report of Project LANA uses for the obstacle at this step
\cite[\S0.2]{LANA}. In each of the
readings listed in \cite{RF}, Corollary~3.12 and IUT~IV together give the same height inequality.

\tagH{} The question of this note is the following. In function fields, an inequality of the same shape follows
structurally from the Kodaira--Spencer morphism. Can its ``number-field version'' be realized, as a quantity, on the
construction of IUT?

\subsection{What this note is}
\begin{itemize}
\item \textbf{An essay that proposes a viewpoint.} Proposition~3.1 is classical, and Propositions~4.1, 4.2, 5.1, 7.1
and 7.2 are elementary. What this note adds is a dictionary that places three situations in one form, the dichotomy
seen from it, and a concrete check that an additive bridge does not cross the $\Theta$-link (end of \S1.3).
\item \textbf{What is not claimed.} We do not say that IUT is wrong, nor that the abc conjecture is false, nor that the
inequality of Corollary~3.12 fails for any single datum (on the six completely factored data of this note it holds in
the reading of the container \cite{TC}). The undisputed parts of IUT (anabelian reconstruction, the \'etale theta
function, Frobenioids) remain mathematics regardless of this note.
\item \textbf{A ``bridge'' in this note is one brought in from outside.} What \S7 shows is that a global relation lying
outside the $\Theta$-link, such as an additive relation, cannot be made to act on the comparison through the
$\Theta$-link. We do not exclude that a bridge of another form is built inside the multiradial representation (\S8.2).
\end{itemize}

\subsection{Prior work}
\begin{itemize}
\item \tagS{} \textbf{The analogy between the Kodaira--Spencer morphism and IUT is Mochizuki's own.} \cite[\S2.9]{Alien}
recalls the classical argument that derives the Szpiro inequality of function fields from the Kodaira--Spencer
morphism, and \cite[\S3.1, (v)]{Alien} reformulates it in terms of multiradiality and calls it analogous to the
``license to confuse''. \cite[\S2.3]{Alien} describes the isogeny argument in the presence of a global multiplicative
subspace and calls the ``simulation'' of such a subspace one fundamental theme of IUT. Mochizuki introduced an
arithmetic Kodaira--Spencer morphism in 1999 \cite[Chapter~IX]{HA}, and in 2000 wrote that its Galois-theoretic version
can only be defined when the Tate module has a natural rank one multiplicative subspace \cite{GalKS}.
\item \tagS{} \textbf{The two requirements on the container --- large enough to house both sides, small enough to give
a nontrivial estimate --- are also stated by Mochizuki.} His slides on the classical roots of IUT give, as the key
point of the Kodaira--Spencer argument, that $p_1^*\Ecal\cong p_2^*\Ecal$ serves as a common container, and say that
one way to understand the multiradial representation of IUT~III is as the construction of a common container that is
\q{sufficiently large to house the incompatible ring structures on either side of the gluing constituted by the theta
link $q_E^N\mapsto q_E$}, but \q{sufficiently small to yield nontrivial estimate on the height of the elliptic curve
over a number field under consideration} \cite[\S2]{Roots}. In \cite[\S3.10, (ii), (iii)]{EssLgc} he writes that
\q{this approach consists of embedding $Y_1$ into larger and larger containers [i.e., the containers obtaining by
adding the various indeterminacies] until the container becomes sufficiently large as to contain/engulf [not only
$Y_1$, but also (!)] $Y_2$}, and compares this with the PD-envelopes of the theory of crystals \cite{EssLgc}. The
``bridge'' and the ``size'' of this note (Definition~2.1) restate these two requirements in the language of the model
of \cite{SL}.
\item \tagX{} \textbf{That an arithmetic Kodaira--Spencer map is not yet available} is also stated in Fesenko's survey:
\q{A (full) arithmetic version of the Kodaira--Spencer map could be quite useful for giving a proof of the arithmetic
Szpiro conjecture. However, such an arithmetic version of the Kodaira--Spencer map is not yet known.}
\cite[\S1.3]{Fe}. In \S1.4 he adds that \q{His earlier Hodge-Arakelov theory [23], where a certain weak arithmetic
version of the Kodaira--Spencer map is studied, was already an innovative step forward} \cite{Fe}.
\item \tagX{} \textbf{A Kodaira--Spencer class for number fields does not exist, depending on the definition.} Faltings
showed that a Kodaira--Spencer class for curves over number fields satisfying certain compatibility axioms cannot exist
\cite{Fal}; Dupuy--Zureick-Brown showed that, with slightly modified definitions, the Deligne--Illusie classes satisfy a
compatibility condition of this kind \cite{DZB}.
\item \tagX{} \textbf{The equivalence of arithmetic derivatives adapted to an additive relation with the abc
conjecture} is due to Pasten \cite{Pa}. On $a+b=c$, arithmetic derivatives that are additive exist unconditionally, and
the existence of sufficiently small ones is equivalent to the abc conjecture. Half of the view of this note, ``the
bridge is unconditional, the size is abc'', is this result itself. Laniewski and M\"uller (arXiv:2610.26490, 19 October
2026) place Pasten's inequality in the quality plane of abc triples and show that \q{the largest absolute value of the
derivation on $a$, $b$ and $c$ exceeds $c/\mathrm{rad}(abc)$ whenever the derivation is additive on the triple with
non-zero Wronskian, so that it gives no bound on the quality by itself} \cite{LM}.
\item \tagX{} \textbf{A weak abc from linear forms in logarithms} is due to Stewart--Yu \cite{StYu}; \cite[\S3.7,
(iii)]{Alien} compares it with IUT.
\item \tagX{} \textbf{The general theory of height inequalities and the Kodaira--Spencer map for function fields} is in
M.~Kim \cite{Kim}.
\item \tagX{} \textbf{In Lean}, lana-agents/iut proves the Szpiro inequality with exponent $6$ over $k[X]$ from the
Mason--Stothers theorem, and that the derivative hypothesis cannot be dropped in characteristic $p$, where Frobenius
twists violate every Szpiro-type bound (\lean{Iut/\allowbreak{}Abc/\allowbreak{}FunctionFieldSzpiro.lean}, 14 October
2026); and it checks by \lean{decide} a table of the ingredients that a proof of the Szpiro inequality over $\Q$ must
use, one of which is the IUT coupling of the wall (\lean{Iut/\allowbreak{}Abc/\allowbreak{}P24Table.lean}, 14 October
2026) \cite{lana}. Both are related in topic to \S\S3--4.
\item \tagX{} \textbf{Joshi} states that a paper treating Kodaira--Spencer theory in his theory of arithmetic
Teichm\"uller spaces (Construction~V) is in preparation \cite{JoQA}. This note does not know its contents and does not
use the outline in the questions and answers.
\item \tagX{} \textbf{The word ``bridge'' was used earlier for the same inclusion.} Higuchi calls it the
\q{numerical bridge statement} and proves that \q{the bridge statement is not a logical consequence} of a qualitative
package extracted from the passage (31 March 2026) \cite{Hi}. Pajo proves a no-go theorem for \q{a deliberately modest
typed abstraction of a common-container-to-numerical bridge} that excludes any \q{bridge certificate} (26 April 2026)
\cite{Pj}. Geiger's \emph{bridge certificate} \q{fixes seven obligations that any admissible passage from a qualitative
linkage to a numerical membership statement must discharge} and is applied to the proof of Corollary~3.12 (13--14 August
2026) \cite{Ge3}; his \emph{bridge contract} follows (16 August 2026) \cite{Ge2}. Definition~\ref{d:bridge} uses the word
for the existence of a lax 2-cell in the model of \cite{SL}, which for IUT is this inclusion.
\item \tagX{} \textbf{That Step~(xi) does not follow from the written properties} is shown in lana-agents/iut
\cite{lana} (\lean{NonDerivable.lean}), by Higuchi \cite{Hi} and in \cite[Theorem~5.7]{SL}. Scholze--Stix write of
Mochizuki's claim that \q{up to the `blurring' given by certain indeterminacies the diagram does commute; it seems to us
that this statement means that the blurring must be by a factor of at least $O(\ell^2)$ rendering the inequality
thus obtained useless} \cite[\S2.2]{SS}.
\item \tagH{} \textbf{What this note adds} is limited to three things. (a) A dictionary that places the Kodaira--Spencer
morphism of function fields, the bridges from an additive relation (Baker, Pasten) and Step~(xi) of IUT in the single
form of a lax 2-cell of \cite{SL}. (b) The assignment of the three situations by ``bridge'' and ``size'' (function
fields have both, Baker and Pasten only the bridge, IUT only a bound on the blurring). The two requirements themselves
are Mochizuki's (the item above); what this note adds is to state them as the existence of a lax 2-cell and its minimal
blurring, and to put the additive bridges (Baker, Pasten) into the same table. For IUT the assignment is earlier.
Geiger's audit of the passage (16 August 2026) finds the load of the theta values quantified in IV, \q{carried through
that computation to the closed form $(\ell+1)/24$}, and records \q{F-R4a therefore delivers for the composite of the two
volumes; F-R4b does not, that section taking the corollary's conclusion as an input and quantifying downstream of the
bridge} and \q{F-R6b does not deliver, and it is the principal chain-closing obligation} \cite[\S6.2]{Ge2}. The request
in Takkun-kohinata/IUT\_LEAN (15 July 2026) states, in its observation (F4), that the computation of IV is correct and
verified, while the derivation that justifies reading it as the containment of Step~(v) is not (paraphrased from the
Japanese) \cite{TK}. Version
1.0.0 of this note said that it had not found this assignment in the sources it had checked (Pasten,
Dupuy--Zureick-Brown, Fesenko, Joshi's questions and answers, Mochizuki's papers and slides, lana-agents/iut, and a web
search on 6 November 2026); Geiger's papers and Takkun-kohinata/IUT\_LEAN were not among them. The reading points to the
same place as Scholze--Stix, lana-agents/iut and Higuchi. (c) A check, along the definitions, that an additive
bridge does not cross the $\Theta$-link, with numbers on actual data (\S7). Item (c) is a concrete instance of the
point of Scholze--Stix and of the incompatibility of the $\Theta$-link with the ring structure (arithmetic degrees)
that Mochizuki himself describes (\cite[\S14, (LbEx1)]{R18}, \cite[\S3.3, (iv)]{Alien}); it is not new in principle.
\end{itemize}

\subsection{The shape of the argument}
\begin{enumerate}
\item \textbf{Separate the bridge from the size} (\S2). We recall the lax 2-cells of \cite{SL} and define ``bridge'' and
``size''.
\item \textbf{In function fields both are automatic} (\S3). The Kodaira--Spencer morphism is a lax 2-cell that exists
automatically by the structure, with blurring $\deg\omega(\log)$.
\item \textbf{A bridge from an additive relation is unconditional; its size is abc} (\S4). The $p$-adic logarithm
bridge (Baker, Stewart--Yu) and the arithmetic derivatives of Pasten.
\item \textbf{The isogeny bridge lies outside the domain of IUT} (\S5). Global multiplicative subspaces and
Definition~3.1~(c).
\item \textbf{In IUT the blurring is small, but the bridge is not derived} (\S6). The statement of Theorem~3.11, where
the source says the $\wedge$ is
derived, and where elements specific to the curve enter.
\item \textbf{An additive bridge does not cross the $\Theta$-link} (\S7). What a full poly-isomorphism carries, the
dilemma at the level of units, the language of HoTT, and what HoTT contributed.
\item \textbf{The dichotomy, and what a repair needs} (\S8).
\end{enumerate}

\subsection{Labels}
\tagS{} a statement of the sources (IUT and Mochizuki's papers). \tagX{} a claim of others (Pasten, Faltings,
Dupuy--Zureick-Brown, Fesenko, lana-agents/iut, Scholze--Stix and others). \tagP{} a statement proved in this note, or a
fact shown in the earlier notes; a statement with only \tagP{} is elementary. \tagR{} the mathematics is known, and the
way it is applied is this note's; the ``known mathematics'' is indicated. \tagH{} a reading, interpretation or proposal
of this note (not a proof).

\section{Bridge and size: lax 2-cells}
\tagP{} We briefly recall the model of \cite[\S5]{SL}. The value of the same object $b$ is carried to the real numbers
along two paths. One is the path $\gamma$ through the $\Theta$-link, which gives $\gamma_*b=-Q$ (the value on the $q$
side). The other is the path $c$ through the shared data, which gives $c_*b=-kQ$ (the value on the $\Theta$ side). The
monodromy of the loop formed by the two paths is $1/k\ne1$. (\cite{SL} writes $\ell$ for the path $\gamma$ and
$m=1/k$ for the scale of the monodromy; we write $\gamma$ to avoid a clash with the prime $\ell$.)
\begin{itemize}
\item A strict 2-cell ($\gamma_*b=c_*b$) does not exist if $Q>0$ \cite[Theorem~5.2]{SL}.
\item Corollary~3.12 (in the model) is equivalent to the lax 2-cell $\gamma_*b\le c_*b+E$, which is equivalent to
$(k-1)Q\le E$ \cite[Theorem~5.3]{SL}.
\item If IUT~IV bounds the blurring by $E\le C$ ($C$ independent of $Q$), then Corollary~3.12 gives $Q\le C/(k-1)$
\cite[Corollary~5.4]{SL}.
\item Qualitative properties such as IPL, SHE and APT (properties of 1-cells) do not decide whether a lax 2-cell exists
\cite[Theorem~5.7]{SL}.
\end{itemize}

\begin{definition}[Bridge and size]\label{d:bridge}
\tagH{} In a situation where values carried along two paths are compared, we call the assertion that a lax 2-cell
exists for some blurring $E$ a \emph{bridge}. When a bridge exists, we call the smallest blurring that makes it possible
its \emph{size}. A strict 2-cell is a bridge of size $0$.
\end{definition}

\tagX{} The word ``bridge'' was used earlier for the same inclusion of Step~(xi), by Higuchi (31 March 2026)
\cite{Hi}, Pajo (26 April 2026) \cite{Pj} and Geiger (August 2026) \cite{Ge3,Ge2}; see \S1.3.

\tagH{} Definition~\ref{d:bridge} restates, in the language of the model of \cite{SL}, the two requirements that
Mochizuki states for the Kodaira--Spencer morphism and for the multiradial representation (that the common container be
large enough to house both sides, and small enough to give a nontrivial estimate; \S1.3). The bridge corresponds to the
former and the size to the latter.

\tagH{} To obtain the height inequality from Corollary~3.12, one needs both a bridge and an upper bound $E\le C$ on its
size. As we shall see, the three situations differ in which of the two they have.

\section{Function fields: the Kodaira--Spencer morphism is a canonical lax 2-cell}
\tagS{} \cite[\S2.9]{Alien} constructs the Kodaira--Spencer morphism, for a family of elliptic curves over a log curve
$S^{\log}$ over $\mathbb C$, from the de Rham bundle $\Ecal$, the Hodge subbundle $\omega_E\subseteq\Ecal$ and the
Gauss--Manin connection $\nabla_{\Ecal}$, and reads it as follows.
\srcquote{\tagS{} \q{this diagram may be understood as a computation of the degree $\deg(-)$ of the `positive mass'
Frobenius-like object $\omega_E$ that yields an inequality} \dots{} \q{via a comparison of `alien copies' of the
`positive mass' Frobenius-like object $\omega_E$ that lie on opposite sides of the `link' constituted by a
deformation of the moduli/holomorphic structure of the family of elliptic curves under consideration}}{Mochizuki,
\emph{The mathematics of mutually alien copies}, \S2.9 \cite{Alien}}
Here the inequality is $2\cdot\deg(\omega_E)\le\deg(\omega_{S^{\log}/\mathbb C})$. \cite[\S3.1, (v)]{Alien} restates the
same argument in five steps $(1^{\KS})$--$(5^{\KS})$. One takes a line bundle $\Lcal$ and a rank two bundle $\Fcal$
with an exact sequence $0\to\Lcal\to\Fcal\to\Lcal^{-1}\to0$. On the first infinitesimal neighbourhood
$S^{\log}_\delta$ of the log diagonal, the isomorphism $p_1^*\Fcal\simeq p_2^*\Fcal$ of the two pull-backs given by the
connection lets one compare the two pull-backs $p_1^*\Lcal$ and $p_2^*\Lcal$, with $\Fcal$ as a ``container'' for
$\Lcal$. If the two images coincide, then $\deg\Lcal=0$ ($(4^{\KS})$); if not, the discrepancy gives a nonzero morphism
of line bundles $\Lcal\to\omega_{S^{\log}/\mathbb C}\otimes\Lcal^{-1}$ ($(5^{\KS})$). For $\Fcal=\Ecal$ and
$\Lcal=\omega_E$ this is the Kodaira--Spencer morphism.
\srcquote{\tagS{} \q{the Kodaira-Spencer morphism may be thought of as a measure of the discrepancy that arises when one
fixes the `$\omega_E$' on one factor of $S$ and compares it with the `$\omega_E$' on a distinct, `alien'
factor of $S$ by means of the common container `$\Ecal$', which is equipped with a connection
`$\nabla_{\Ecal}$'} \dots{} \q{When formulated in this way, the Kodaira-Spencer morphism becomes manifestly analogous to
the approach sketched in the Introduction to the present paper} \dots{} \q{i.e., [in the language of the
Introduction] a `license to confuse'}}{ibid., \S3.1, (v) \cite{Alien}}
\srcquote{\tagS{} \q{in the limit `$N\to\infty$' [in which case comparison between `$h$' and `$N\cdot h$',
or equivalently, between `$h$' and `$\frac1N\cdot h$', becomes a comparison between `$h$' and `$0$']}}%
{ibid., \S3.1, (v) (on the relation between the classical argument and the argument of IUT that multiplies the height
$h$ by $N$) \cite{Alien}}

\begin{proposition}[The Kodaira--Spencer morphism is a lax 2-cell]\label{p:ks}
\tagR{}\known{Known mathematics: classical (\cite[\S2.9, \S3.1, (v)]{Alien}, \cite{Kim}). The way it is applied: the
2-cells of \cite{SL}.}
Let $S$ be a proper smooth curve over $\mathbb C$, $D\subseteq S$ a reduced divisor and $\omega=\Omega_S(\log D)$. Let
$\Fcal$ be a vector bundle of rank $2$, $\nabla:\Fcal\to\Fcal\otimes\omega$ a logarithmic connection whose residues are
nilpotent at every point, and $\Lcal\subseteq\Fcal$ a line subbundle with $\Fcal/\Lcal\simeq\Lcal^{-1}$. For the
composite
\[\kappa:\Lcal\hookrightarrow\Fcal\xrightarrow{\nabla}\Fcal\otimes\omega\to(\Fcal/\Lcal)\otimes\omega
=\Lcal^{-1}\otimes\omega\]
we have:
\begin{enumerate}[label=(\roman*)]
\item $\kappa$ is $\Ocal_S$-linear;
\item if $\kappa=0$, then $\deg\Lcal=0$ (a strict 2-cell);
\item if $\kappa\ne0$, then $2\deg\Lcal\le\deg\omega=2g-2+\#D$ (a lax 2-cell with blurring $\deg\omega$).
\end{enumerate}
\end{proposition}
\begin{proof}
(i) For $f\in\Ocal_S$ and $s\in\Lcal$, $\nabla(fs)=f\nabla s+s\otimes df$, and $s\otimes df$ lies in
$\Lcal\otimes\omega$, hence vanishes in $(\Fcal/\Lcal)\otimes\omega$. (ii) If $\kappa=0$, then $\nabla$ preserves $\Lcal$
and gives a logarithmic connection on $\Lcal$. Its residues are restrictions of nilpotent residues to an invariant line,
hence $0$, and the residue theorem gives $\deg\Lcal=-\sum\mathrm{res}=0$. (iii) There is a nonzero morphism of line
bundles $\Lcal\to\Lcal^{-1}\otimes\omega$, so $\deg\Lcal\le\deg\omega-\deg\Lcal$.
\end{proof}

\tagH{} In the language of \cite{SL}, the two pull-backs are the two paths, the connection is the shared data (the
common container), the coincidence of the images is a strict 2-cell, and $\kappa$ is a lax 2-cell. \textbf{Both the
bridge and the size come automatically from the structure.} The bridge comes from the fact that $\kappa$ is a morphism
that is linear over the base ((i)). The size $\deg\omega$ is determined by the curve alone and does not depend on the
height $\deg\Lcal$ of the family. This linearity comes from the fact that the comparison takes place over a single base
$S$ and its ring structure, and moreover in an infinitesimal neighbourhood. Number fields have neither derivations of
this kind nor infinitesimal neighbourhoods (\S4.5).

\section{Number fields (1): bridges from an additive relation}

\subsection{The \texorpdfstring{$p$}{p}-adic logarithm bridge}
\tagH{} In number fields, the closest thing to the discrepancy of Kodaira--Spencer is to linearize an additive relation
by the $p$-adic logarithm. The exponent of a prime factor of one term (information in the value group) moves to the
logarithm of the ratio of the other two terms (information in the unit group). In the words of Remark~3.12.2~(ii) of
IUT~III, this amounts to computing the value group portion in terms of the unit group portion; here this is done by an
addition rather than by theta functions.
\srcquote{\tagS{} \q{this `reassembling procedure' allows one to compute/estimate the value group portions of various
monoids of arithmetic interest in terms of the unit group portions of these monoids. It is precisely these estimates
that give rise to the inequality obtained in Corollary 3.12.}}{IUT~III, Remark~3.12.2, (ii) \cite{IUT3}}

\begin{proposition}[The additive bridge]\label{p:additive}
\tagP{}\known{Elementary. Numerical check: \texttt{scripts/bridge\_lib.py}.}
Let $a,b,c$ be coprime positive integers with $a+b=c$, $p$ an odd prime and $e=v_p(a)\ge1$. Then:
\begin{enumerate}[label=(\roman*)]
\item $b/c\in1+p^e\Z_p$ and $v_p(b/c-1)=e$;
\item $v_p(\log_p(b/c))=e$;
\item $\log_p(b/c)=\sum_{q\mid bc}\bigl(v_q(b)-v_q(c)\bigr)\log_pq$, that is, $\log_p(b/c)$ is a linear form in the
$p$-adic logarithms of the primes dividing $bc$, with the exponents as coefficients.
\end{enumerate}
The same holds for $a/c$ if $p\mid b$, and for $-a/b$ if $p\mid c$.
\end{proposition}
\begin{proof}
(i) $b=c-a\equiv c\pmod{p^e}$ and $p\nmid c$ (coprimality), so $b/c-1=-a/c$ has valuation exactly $e$. (ii) If $p$ is
odd and $v_p(x)=e\ge1$, then for $n\ge2$ we have $v_p(x^n/n)\ge ne-v_p(n)>e$ (as $v_p(n)\le\log_3n<n-1$), so
$\log(1+x)=x-x^2/2+\cdots$ has valuation $e$. (iii) $\log_p$ is a homomorphism on $\Z_p^\times$, and every prime
dividing $bc$ is different from $p$.
\end{proof}

\tagP{} We checked this on well-known abc triples of high quality, computing $\log_p$ from the series
(Table~\ref{t:triples}, which lists all odd primes with exponent at least $2$).

\begin{table}[ht]\centering\small
\caption{The additive bridge: the exponent $e$ in one term and the valuation of the $p$-adic logarithm of the ratio of
the other two.}\label{t:triples}
\begin{tabular}{@{}llllllll@{}}
\toprule
triple $a+b=c$ & quality & $p^e$ & term & ratio & $v_p(\log_p(\text{ratio}))$ & primes in the ratio\\
\midrule
\tabletriples
\bottomrule
\end{tabular}
\end{table}

\tagH{} For example, in $2+3^{10}\cdot109=23^5$ the high power $3^{10}$ moves to the fact that the linear form
$\log_32-5\log_323$ in the $3$-adic logarithms of $2$ and $23$ is unusually small $3$-adically (valuation $10$). The
coefficients of the linear form are the exponents themselves ($1$ and $5$), which are small.

\subsection{Initial \texorpdfstring{$\Theta$}{Theta}-data in Legendre form}
\tagS{} In the application to ABC, IUT~IV uses the line of parameters of elliptic curves in Legendre form.
\srcquote{\tagS{} \q{We shall regard $X$ as the `$\lambda$-line'}}{IUT~IV, Corollary~2.2 \cite{IUT4}}

\begin{proposition}[The depth of the bridge on initial $\Theta$-data]\label{p:legendre}
\tagP{}\known{Elementary.}
Let $F$ be a number field, $\lambda\in F\setminus\{0,1\}$ and $E_\lambda:y^2=x(x-1)(x-\lambda)$. Let $v$ be a place of
odd residue characteristic at which $E_\lambda$ has (potentially) multiplicative reduction. In each of the cases
$\mathrm{ord}_v(\lambda)=e>0$, $\mathrm{ord}_v(1-\lambda)=e>0$ and $\mathrm{ord}_v(\lambda)=-e<0$, the element
$1-\lambda$, $\lambda$, $\lambda/(\lambda-1)$ respectively is congruent to $1$ modulo the $e$-th power of the prime
ideal of $v$, and $\mathrm{ord}_v(q)=2e$. Since these three cases are the only ones with multiplicative reduction at an
odd place, at every odd bad place the depth of the complementary unit is $d_v=\mathrm{ord}_v(q)/2$.
\end{proposition}
\begin{proof}
$j(E_\lambda)=2^8(\lambda^2-\lambda+1)^3/(\lambda^2(\lambda-1)^2)$, and at a place of multiplicative reduction
$\mathrm{ord}_v(q)=-\mathrm{ord}_v(j)$. In the three cases one counts $\mathrm{ord}_v(j)=-2e$ directly ($v$ is odd, so
$2^8$ is a unit). The congruences follow from $1-\lambda=1-(\text{a small element})$ and so on.
\end{proof}

\tagP{} The seven initial $\Theta$-data of this note are all in Legendre form over quadratic fields (the \texttt{data/}
of \cite{TC}). D211 has $F=\Q(\sqrt5)$ and $\lambda=\pi^{29}$ ($\pi=16+3\sqrt5$, of norm $211$); at the place $211+$
above $211$, $\mathrm{ord}(q)=\dordq$, that is, $1-\lambda\equiv1\pmod{\pi^{29}}$, of depth $\ddepth$. The norm of
$1-\lambda$ is $4\cdot3^2\cdot5\cdot59\cdot80621$ times two primes of $21$ and $39$ digits, and at the odd places
dividing these the depth is $1$ (the places above $2$ are not bad places).

\subsection{Baker's method and Stewart--Yu}
\tagX{} The smallness in Proposition~\ref{p:additive} is the starting point of the method of linear forms in
logarithms. Lower bounds for how $p$-adically small a linear form $\Lambda=\sum n_q\log_pq$ can be (the $p$-adic version
of Baker's method) bound the exponent $e$ from above in terms of the primes dividing $bc$ and the size of the exponents.
Stewart--Yu used estimates of this kind to prove a weak form of abc \cite{StYu}. We have not checked their text; the
shape of the conclusion is taken from the summary in \cite{Alien}.
\srcquote{\tagS{} \q{Recall that in [StYu], an inequality is obtained which may be thought of as a sort of `weak version
of the ABC Conjecture', i.e., which is, roughly speaking, weaker than the inequality of the usual ABC Conjecture in that
it contains an undesired exponential operation `$\exp(-)$' in its upper bound.}}{Mochizuki, \emph{The mathematics
of mutually alien copies}, \S3.7, (iii) \cite{Alien}}
\tagS{} Of this exponential discrepancy Mochizuki writes that \q{This sort of deviation from the inequality of the usual
ABC Conjecture is of interest from the point of view of the `vertical shift' discussed in \S3.6, (iv)}, and he
identifies the vertical shift with a single iterate of the log-link in the log-theta-lattice (ibid.) \cite{Alien}.

\tagH{} Here \textbf{the bridge is automatic}: $b/c\equiv1\pmod{p^e}$ follows from $a+b=c$ itself. Moreover, this bridge
holds only on actual triples and distinguishes the curves (triples). What is weak is the \textbf{size}, which is
determined by an estimate of transcendence theory, the lower bound for linear forms.

\subsection{Pasten's arithmetic derivatives}
\tagX{} Pasten wrote this shape in the language of derivatives \cite{Pa}. One takes the free module $\Omega$ generated
by variables $\xi_p$, one for each prime, and the map $d:\Z\to\Omega$, $dn=n\sum_{p\mid n}(v_p(n)/p)\,\xi_p$, which
satisfies the Leibniz rule. A group homomorphism $\psi:\Omega\to\Z$ with finite norm $\|\psi\|=\sup_p|\psi(\xi_p)|$ is
called a derivation, and $d^\psi=\psi\circ d$. Let $\Tcal(a,b)$ be the group of derivations with $\psi(\xi_p)=0$ for
$p\nmid abc$ and $d^\psi a+d^\psi b=d^\psi c$. If the Wronskian $W^\psi(a,b)=a\,d^\psi b-b\,d^\psi a$ is nonzero, $a$
and $b$ are called $\psi$-independent.
\begin{itemize}
\item \tagX{} $\Tcal(a,b)$ has rank $r=\omega(abc)-1$, and by Siegel's lemma there are independent
$\psi_1,\dots,\psi_r$ with $\prod_{i=1}^r\|\psi_i\|\le\frac{\omega(abc)}{2\log2}\,c\log c$ (Theorem~2.6). For at least
one of them $a,b$ are $\psi$-independent, and this $\psi$ satisfies $\|\psi\|\le\frac{\omega(abc)}{2\log2}\,c\log c$
(Lemma~3.5).
\item \tagX{} If $a,b$ are $\psi$-independent, then $\frac{c}{\log c}\le\mathrm{rad}(abc)\cdot\frac{\|\psi\|}{\log2}$
(Theorem~3.3; the same shape as Snyder's proof of Mason's theorem).
\item \tagX{} The small derivatives conjecture (Conjecture~3.9): for some $0<\eta<1$, for all triples except finitely
many and those of the form $(1,N,q)$ with $q$ prime (up to order), there is $\psi\in\Tcal(a,b)$ with $\|\psi\|<c^\eta$ that makes $a,b$
$\psi$-independent. It is equivalent to the abc conjecture, with a correspondence of exponents (Lemma~4.1,
Theorem~4.5, Corollary~4.6; precisely, it follows from the abc conjecture of Masser--Oesterl\'e and implies the abc
conjecture of Oesterl\'e with some exponent).
\end{itemize}
\srcquote{\tagX{} \q{the arithmetic analogue of the Mason-Stothers theorem is the abc Conjecture, but the proof for
polynomials heavily uses derivatives and it is unclear how to adapt it to $\Z$.} \q{In fact, we show that the existence
of sufficiently small arithmetic derivatives in our sense is equivalent to the abc Conjecture.}}{Pasten, \S1.1 and
abstract \cite{Pa}}
\tagH{} In the language of 2-cells, $W^\psi(a,b)\ne0$ is the bridge and $\|\psi\|$ is the size. \textbf{The bridge
exists unconditionally (Lemma~3.5), and bounding its size is equivalent to the abc conjecture.} Putting the size
$\frac{\omega}{2\log2}c\log c$ of Lemma~3.5 into Theorem~3.3 gives only a trivial inequality; Pasten writes that
\q{Lemma 3.5 combined with Theorem 3.3 falls short of proving the abc Conjecture} \cite{Pa}, and Laniewski and
M\"uller show that the largest value of such a derivation on $a$, $b$ and $c$ exceeds $c/\mathrm{rad}(abc)$, so that it
gives no bound on the quality by itself \cite{LM} (\S1.3). The derivative $d/dt$ of a
function field lowers degrees by exactly $1$, that is, it has ``size $1$'', which is why the Mason--Stothers inequality
is sharp.

\subsection{Faltings's non-existence result and Dupuy--Zureick-Brown}
\srcquote{\tagX{} \q{We show that an analog of the Kodaira--Spencer class for curves over number-fields cannot exist.}
\q{if $S$ is the spectrum of the ring of integers of a number field then there are no derivations and hence the
Kodaira--Spencer map doesn't make sense.} \q{By modifying his definitions slightly, we show that the Deligne--Illusie
classes satisfy what could be considered an `arithmetic Kodaira--Spencer' compatibility condition.}}%
{Dupuy--Zureick-Brown, introduction (the first sentence is quoted there as the abstract of Faltings's paper)
\cite{DZB}}
\tagH{} The linearity of the function field argument (Proposition~\ref{p:ks}~(i)) came from derivations over the base.
The ring of integers of a number field has none. Pasten avoids this by building additive ``derivations'' only on a
single equation $a+b=c$, and in exchange loses control of the size. The Deligne--Illusie classes give another way
around from the side of Frobenius lifts. In the language of this note, both are other devices for building a bridge.

\section{Number fields (2): global multiplicative subspaces and isogenies}
\tagH{} There is a special case in which both a bridge and a size are available in number fields. \tagS{} It is the
case in which the $\ell$-torsion points of the elliptic curve have a Galois-stable subspace of rank one (a global
multiplicative subspace).
\srcquote{\tagS{} \q{The quotient of the original given elliptic curve by such a global multiplicative subspace is an
elliptic curve that is isogenous to the original elliptic curve. Moreover, the $q$-parameters of this isogenous
elliptic curve are the $l$-th powers of the $q$-parameters of the original elliptic curve} \dots{} \q{On the
other hand, in general, such a global multiplicative subspace does not exist, and the issue of somehow `simulating' the
existence of a global multiplicative subspace is one fundamental theme of inter-universal Teichm\"uller theory.}}%
{Mochizuki, \emph{The mathematics of mutually alien copies}, \S2.3 \cite{Alien}}
\tagS{} According to the same place, the effect of the isogeny on differentials ($d\log(q)\mapsto l\cdot d\log(q)$)
adds, at the level of heights, terms of the order of $\log(l)$, \q{thus giving rise to inequalities that are roughly of
the form `$l\cdot h\le h+\log(l)$'}, and \q{this is precisely the approach that is taken to bounding heights in the
`tiny' special case of the theory of [Falt1] that is given in the proof of [GenEll], Lemma 3.5} \cite{Alien} (there, [Falt1] is
Faltings's paper of 1983 on finiteness theorems; see also \cite{GenEll}). \tagH{} Here the bridge is an algebraic map (compatible with the ring
structure), the isogeny, and the size is of the order of $\log l$. \cite[\S3.1, (v)]{Alien} also applies $(2^{\KS})$ and
$(4^{\KS})$ of the Kodaira--Spencer argument to this situation.
\srcquote{\tagS{} \q{relative to the analogy with inter-universal Teichm\"uller theory, the situation that arises in
$(2^{\KS})$, $(4^{\KS})$ corresponds to the [unusual!] situation in which there actually exists a `global
multiplicative subspace'}}{ibid., \S3.1, (v) \cite{Alien}}
\tagS{} Mochizuki's Galois-theoretic arithmetic Kodaira--Spencer morphism of 2000 can also be defined only where a
multiplicative subspace exists.
\srcquote{\tagS{} \q{Unfortunately, however, this Lagrangian arithmetic Kodaira-Spencer morphism, which is based on a
certain `Lagrangian Galois action,' can only be defined when there is a natural (rank one) multiplicative subspace
(i.e., `weight 1' subspace) of the Tate module of the elliptic curve in question.} \dots{} \q{Ultimately, however, one
would like to carry out this construction for elliptic curves over number fields.}}{Mochizuki, \emph{The
Galois-theoretic Kodaira-Spencer morphism of an elliptic curve}, \S0 \cite{GalKS}}
\tagS{} IUT~I regards the whole theory as a solution of the problem of constructing such a subspace, and states the
price of ``simulating'' it.
\srcquote{\tagS{} \q{a sort of solution to the problem of constructing `global multiplicative subspaces' and `global
canonical generators'} \dots{} \q{the nonexistence of which in a `naive, scheme-theoretic sense' constitutes the main
obstruction to applying the theory of [HASurI], [HASurII] to diophantine geometry} \dots{} \q{one can, in some sense,
`simulate' the notions of a `global multiplicative subspace' or a `global canonical generator'. On the other hand, such
`simulated global objects' may only be achieved at the cost of `dismantling', or performing `surgery' on, the global
prime structure of the number fields involved}}{IUT~I, introduction \cite{IUT1}}
\srcquote{\tagS{} \q{this sort of operation of discarding certain of the primes of an NF can only be performed if one
forgets the additive structure of an NF [i.e., since a sum of elements of an NF that are invertible at a given nonempty
set of primes is no longer necessarily invertible at those primes!} \dots{} \q{and works only with multiplicative
structures, e.g., with Frobenioids.}}{Mochizuki, \emph{The mathematics of mutually alien copies}, \S3.3, (iv)
\cite{Alien}}
\srcquote{\tagS{} \q{$l$ is a prime number $\ge5$ such that the image of the outer homomorphism
$G_F\to\mathrm{GL}_2(\mathbb F_l)$ determined by the $l$-torsion points of $E_F$ contains the subgroup
$\mathrm{SL}_2(\mathbb F_l)\subseteq\mathrm{GL}_2(\mathbb F_l)$}}{IUT~I, Definition~3.1, (c) \cite{IUT1}}

\begin{proposition}[The isogeny bridge is excluded from the initial $\Theta$-data]\label{p:sl2}
\tagP{}\known{Elementary.}
If the condition of Definition~3.1~(c) holds, then $E_F[\ell]$ has no $G_F$-stable line. In particular, there is no
global multiplicative subspace for this $\ell$.
\end{proposition}
\begin{proof}
$\mathrm{SL}_2(\mathbb F_\ell)$ acts transitively on the nonzero vectors of $\mathbb F_\ell^2$, so it fixes no line.
\end{proof}

\tagH{} The curves on which the isogeny bridge acts are excluded from the domain of IUT by definition. IUT deliberately
treats the general case, in which no bridge of this kind is available. Moreover, in the application to ABC, $\ell$ is
taken of the order of the square root of the height (IUT~IV, Remark~2.2.4); even if a curve has an isogeny for some
small $\ell$, that $\ell$ cannot be used.
\srcquote{\tagS{} \q{The prime number $l$ is, ultimately, in the computations of Corollary 2.2, (ii) [cf.,
especially, condition `(C1)'], taken to be roughly of the order of the square root of the height of the elliptic curve
under consideration.}}{IUT~IV, Remark~2.2.4 \cite{IUT4}}

\section{IUT: the blurring is small, but the bridge is not derived}

\subsection{The statement of Theorem 3.11}
\tagS{} Theorem~3.11 only fixes initial $\Theta$-data as in Definition~3.1 of IUT~I and adds no condition.
Compatibility with the $\Theta$-link is defined as stability, equivariance and functoriality with respect to
automorphisms.
\srcquote{\tagS{} \q{Fix a collection of initial $\Theta$-data} \dots{} \q{as in [IUTchI], Definition 3.1.}
\dots{} \q{whose associated `global degrees' may be computed by means of the log-volumes of (a)} \dots{} \q{in the sense
that these constructions are stabilized/equivariant/functorial with respect to arbitrary automomorphisms}
[\emph{sic}] \q{of the domain and codomain of these horizontal arrows of the LGP-Gaussian log-theta-lattice.}}%
{IUT~III, Theorem~3.11, opening, (i)(c), (iii)(c) \cite{IUT3}}
\tagP{} The statement of Theorem~3.11 does not contain the $q$-pilot. The $q$-pilot appears in Remark~3.11.1~(iii),
immediately after, as the property (IPL) that the output is constructed in such a way that it is \q{linked/related}
via full poly-isomorphisms to the $q$-pilot \cite{IUT3}.
\srcquote{\tagS{} \q{This output data is constructed in such a way that it is linked/related, via full
poly-isomorphisms of $\FV$-prime-strips induced by operations in the algorithm, to the input data prime-strip, i.e.,
the `coric'/`fixed' $q$-pilot $\FV$-prime-strip, equipped with its rigidifying arithmetic holomorphic
structure}}{IUT~III, Remark~3.11.1, (iii) (IPL) \cite{IUT3}}
\tagP{} In Step~(xi-f) the inclusion enters in the form that the construction of the subset of possible log-volumes of
the output \q{is subject to the condition that this construction of output data possibilities constitutes, in
particular, a construction [perhaps only up to some sort of `approximation', as a result of various indeterminacies] of
the pilot-object log-volume of the input data} \cite{IUT3}, namely of $-|\log(\qq)|$; it does not follow from the
written properties (IPL, SHE, APT) (\cite[Theorem~5.7]{SL}, \cite{lana}, \cite{Hi}); for the log-volume content of the
22 results cited in the proof, on actual local objects, see \cite[Theorem~10, Corollary~11, Theorem~22]{W311}. On the
other hand, the upper bound
$E\le C$ on the blurring is given by the computation of IUT~IV (the proof of Theorem~1.10), and the dispute lies on the
side of Step~(xi) of Corollary~3.12. \textbf{The written argument gives a small blurring but does not derive the
bridge.}

\subsection{Where the source says the \texorpdfstring{$\wedge$}{AND} is derived}
\tagS{} The source itself names where the $\wedge$ of the comparison is derived from: the essentially tautological
$\wedge$ coming from the fact that the two arithmetic holomorphic structures are distinguished ($(\mathrm
c^{\mathrm{itw}})$) is mapped, through the subtle compatibility of the quotient operations of Remark~3.9.5~(viii), to
the $\wedge$ on weak data that does not distinguish the two ($(\mathrm f^{\mathrm{itw}})$).
\srcquote{\tagS{} \q{data that is `sufficiently weak' that one can no longer distinguish between $\Theta$-pilot
and $q$-pilot objects.} \dots{} \q{it is essentially a tautological consequence of the fact that these two
arithmetic holomorphic structures in the domain and codomain of the $\ThetaLGP$-link are distinguished from one
another} \dots{} \q{where the second `$\Longrightarrow$' of the above display is purely formal} \dots{} \q{the
`$\wedge$' of the above display may be regarded as the `image' of, hence, in particular, as a consequence of, the
`$\wedge$' in the display of $(\mathrm c^{\mathrm{itw}})$, via the various (sub)quotient operations discussed in
Remark 3.9.5, (viii), i.e., whose subtle compatibility properties allow one to conclude the `$\wedge$' of the above
display from the `$\wedge$' in the display of $(\mathrm c^{\mathrm{itw}})$}}{IUT~III, Remark~3.12.2, (ii),
$(\mathrm b^{\mathrm{itw}})$, $(\mathrm c^{\mathrm{itw}})$, $(\mathrm f^{\mathrm{itw}})$ \cite{IUT3}}
\srcquote{\tagS{} \q{the nontriviality of these various (sub)quotient operations lies in the very delicate balance
between minimizing the indeterminacies that arise from passing to a quotient, while at the same time ensuring
compatibility with the structures that exist prior to formation of the quotient.} \dots{} \q{forms a closed loop, i.e.,
up to the introduction of the `formal quotient indeterminacies' discussed in (cQ3), (cQ4)} \dots{} \q{a non-closed loop
would yield a situation from which no nontrivial conclusions may be drawn} \dots{} \q{passing to log-volumes [i.e.,
(sQ5)] amounts precisely to forgetting the local Galois groups and Frobenius-like units}}{IUT~III, Remark~3.9.5,
(viii), (ix) \cite{IUT3}}
\tagH{} The closed loop closes at the level of structures that include the units and the local Galois groups of the
prime-strips, before the passage (sQ5) to log-volumes. In the language of this note, the two sides of the ``delicate
balance'' (keeping the indeterminacies small, and keeping compatibility with the original structures) correspond to a
small blurring ($E\le C$) and the bridge (the inclusion). The source says that this balance is maintained, and the written text
states it at the level of structures.

\subsection{Where elements specific to the curve enter}
\tagH{} Information specific to the curve enters the construction as the ``critical points'' of $\kappa$-coric
functions. \tagS{} A critical point is a point lying over a closed point of the compactification $|C_L|^{\mathrm{cpt}}$
of the coarse space $|C_L|$ that arises from a $2$-torsion point of $E_F$; those not over the cusp are called strictly
critical (for the Legendre form, those over $x=0,1,\lambda$ and the cusp $\infty$).
\srcquote{\tagS{} \q{We shall refer to a closed point of the proper smooth curve determined by some finite subextension}
\dots{} \q{as a critical point if it maps to a closed point of the [proper smooth] compactification} \dots{} \q{that
arises from one of the 2-torsion points of $E_F$} \dots{} \q{$f$ restricts to a root of unity at every strictly critical
point of $|C_L|^{\mathrm{cpt}}$} \dots{} \q{restriction to the strictly critical points determines a sort of canonical
splitting up to roots of unity} \dots{} \q{of the set of nonzero constant [i.e., $L$-] multiples of $\kappa$-coric
functions into a direct product, up to roots of unity, of the set of $\kappa$-coric functions and the set of nonzero
elements of $L$} \dots{} \q{since $|C_L|^{\mathrm{cpt}}$ has precisely 4 critical points}}{IUT~I, Remark~3.1.7, (i),
(ii) \cite{IUT1}}
\tagS{} Restricting the Kummer classes of these functions to the decomposition groups of non-critical points, one
recovers the number field $F_{\mathrm{mod}}$ as a field (IUT~I, Example~5.1~(v): \q{together with the field structure
--- and hence, in particular, the topologies determined by the valuations ---} on the union of the monoids with $\{0\}$)
\cite{IUT1}. IUT~II, Corollary~4.8 states the compatibility of the isomorphisms of this Kummer theory with the labels,
the $\mathbb F_l^{\divideontimes}$-symmetry and localization, and its proof is \q{The various assertions of Corollary 4.8
follow immediately from the definitions and the references quoted in the statements of these assertions.}
\cite{IUT2}. In Theorem~3.11~(i)(c), the copies of the number field are placed in the product of log-shells, and their
role is to compute global degrees by log-volumes (the quotation in \S6.1).

\begin{observation}[Elements specific to the curve enter the construction, not the comparison]\label{o:curve}
\tagH{}\known{A reading on the statements of Theorem~3.11, IUT~II Corollary~4.8, IUT~I Example~5.1 and Remark~3.1.7.}
$\lambda$ enters the construction as a tool for recovering the number field $F_{\mathrm{mod}}$ rigidly (the critical
points of $\kappa$-coric functions, the canonical splitting of the constants). The only global information that enters
the comparison (the inequality between log-volumes) is the global degree of arithmetic line bundles of the number field,
that is, the sum of local log-volumes (the product formula). The product formula holds for every element and so does
not distinguish curves. No statement in these theorems makes the additive relation $\lambda+(1-\lambda)=1$ act on the
comparison as a quantity.
\end{observation}
\tagH{} Applying the product formula to $j(E)$ gives only an identity (the contribution of $q$ at the bad places equals
the contribution of the zeros of $j$ plus the archimedean contribution), not an upper bound. To make it an upper bound
one needs to know how the additive connection acts as a quantity, and that is the abc conjecture itself (\S4).

\section{An additive bridge does not cross the \texorpdfstring{$\Theta$}{Theta}-link}

\subsection{What the \texorpdfstring{$\Theta$}{Theta}-link carries}
\tagS{} The $\Theta$-link carries the unit portions by full poly-isomorphisms, and the categories glued together are
multiplicative.
\srcquote{\tagS{} \q{the full poly-isomorphism of $\FU$-prime-strips} \dots{} \q{it is of crucial importance that
the categories that are glued together via the $\ThetaLGP$-link be purely multiplicative in nature}}{IUT~III,
Theorem~3.11, (iii)(a), and Remark~3.9.5, (x) \cite{IUT3}}

\begin{proposition}[What a full poly-isomorphism carries]\label{p:fullpoly}
\tagP{}\known{Elementary.}
Let $X,Y$ be objects of a category, $\Phi=\mathrm{Isom}(X,Y)$ the full poly-isomorphism and $T$ a functor that assigns
a set to each object. For $x\in T(X)$, the set $\{T(\varphi)(x)\mid\varphi\in\Phi\}$ obtained by carrying $x$ along
$\Phi$ is one $\mathrm{Aut}(Y)$-orbit. It is a single point if and only if $x$ is fixed by $\mathrm{Aut}(X)$. For a
property $P$ invariant under isomorphisms, $P(X)\iff P(Y)$.
\end{proposition}
\begin{proof}
For $\varphi,\varphi'\in\Phi$, $\varphi'=(\varphi'\varphi^{-1})\varphi$, and $\varphi'\varphi^{-1}$ runs over all of
$\mathrm{Aut}(Y)$. The last assertion is the definition.
\end{proof}

\tagH{} The material of an additive bridge consists of the units $u_v=(1-\lambda)_v$ (or their logarithms) and of the
consistency that ``they all come, at every bad place, from one and the same global element''. The former are specific
elements of the unit portion, and, for instance, the automorphisms (Ism) of (Ind2) move $\log u_v$ to other elements. By
Proposition~\ref{p:fullpoly}, therefore, what the $\Theta$-link can carry is a property of the type ``there is a unit of
depth $d_v$'', not the element $1-\lambda$. The latter consistency is not contained in the data of the local
$\FU$-prime-strips, and the global portion only has the realified degrees. The source itself writes that these data
are \q{`sufficiently weak' that one can no longer distinguish between $\Theta$-pilot and $q$-pilot objects}
(\S6.2) \cite{IUT3}.

\subsection{The dilemma at the level of units}
\begin{proposition}[Keeping an additive bridge on both sides of the $\Theta$-link]\label{p:gap}
\tagP{}\known{Elementary. Numbers: \texttt{scripts/bridge\_lib.py}.}
Let $\Vbad$ be the same set of bad places as in \cite{TC} (the places with $\mathrm{ord}_v(q)>0$ that are not above $2$,
$3$, $5$ or $\ell$ and at which $\mathrm{ord}_v(q)$ is not divisible by $\ell$), and let $d_v=\mathrm{ord}_v(q)/2$ be
the depth of the additive bridge at $v\in\Vbad$ (Proposition~\ref{p:legendre}). The $\Theta$-link identifies the
$\Theta$-pilot $\qq^{j^2}$ of label $j$ with the $q$-pilot $\qq$ on the other side, so that, on taking $2\ell$-th
powers, $q^{j^2}$ on this side corresponds to the $q$ of the curve on the other side. Suppose that (a) the units of a
bridge of depth $d_v$ are carried across the $\Theta$-link, (b) the carried units are required to be an additive bridge
on the other side as well, that is, units of depth $\frac12\mathrm{ord}_v(q^{j^2})=j^2d_v$ for the $q$ of the other side
(which corresponds to $q^{j^2}$ on this side), and (c) the operation of carrying changes the depth by at most
$\varepsilon_{v,j}$ levels. Then $(j^2-1)\,d_v\le\varepsilon_{v,j}$, and averaging over the labels and summing with the
weights $f_v\log p_v/[F:\Q]$ gives
\[\frac{k-1}{2}\cdot\frac{1}{[F:\Q]}\sum_{v\in\Vbad}\mathrm{ord}_v(q)\,f_v\log p_v
\;\le\;\frac{1}{[F:\Q]}\sum_{v\in\Vbad}\bar\varepsilon_v\,f_v\log p_v,\]
where $\bar\varepsilon_v$ is the average of $\varepsilon_{v,j}$ over the labels. The left-hand side is exactly half of
the discrepancy at the level of value groups (the average difference between $q$ and $q^{j^2}$) measured with the same
weights. In the normalization of Corollary~3.12,
\[Q=|\log\qq|=\frac{1}{2\ell[F:\Q]}\sum_{v\in\Vbad}\mathrm{ord}_v(q)f_v\log p_v,\]
the left-hand side is $\ell\,(k-1)Q$, that is, $\ell$ times the wall $(k-1)Q$.
\end{proposition}
\begin{proof}
By (c) the depth of the carried units differs from $d_v$ by at most $\varepsilon_{v,j}$, and by (b) it must be
$j^2d_v$, so $(j^2-1)d_v\le\varepsilon_{v,j}$. The average of $j^2$ over the labels $j=1,\dots,\ell^*$ is $k$, so the
average of $j^2-1$ is $k-1$. Insert $d_v=\frac12\mathrm{ord}_v(q)$ and sum. The last equality follows from the formula
for $Q$.
\end{proof}

\tagP{} A unit multiple coming from (Ind2) does not change the depth ($\varepsilon_{v,j}=0$). Hence keeping (a) and (b)
leads to a contradiction on every datum with a bad place. This happens on actual initial $\Theta$-data
(Table~\ref{t:data}).

\begin{table}[ht]\centering\footnotesize\setlength{\tabcolsep}{4.5pt}
\caption{The depth of the additive bridge, and the discrepancy when one tries to keep it on both sides.}\label{t:data}
\begin{tabular}{@{}lrlrrrrrr@{}}
\toprule
datum & $\ell$ & \shortstack[l]{place of the\\largest bridge} & $\mathrm{ord}(q)$ & \shortstack[r]{depth\\$d$} &
\shortstack[r]{gap at\\$j=2$} & \shortstack[r]{gap at\\$j=\ell^*$} & \shortstack[r]{sum of\\gaps (LHS)} &
\shortstack[r]{wall\\$(k-1)Q$}\\
\midrule
\tabledata
\bottomrule
\end{tabular}
\par\smallskip\parbox{0.95\textwidth}{\footnotesize ``Place of the largest bridge'': the place of $\Vbad$ at which
$d_vf_v\log p_v$ is largest. The gaps are numbers of levels (differences of depth). The left-hand side and the wall are
in the normalization of Proposition~\ref{p:gap}; their ratio is exactly $\ell$ on every datum. $^{\dagger}$D2357 does
not count the places above an unfactored cofactor, so it is conditional (on this datum the wall exceeds $C$
\cite{TC}).}
\end{table}

\tagH{} Hence there are three ways to repair IUT by bringing in an additive bridge from outside. If the bridge is
kept on both sides, the contradiction of Scholze--Stix appears at the level of units. If it is carried without being
kept, the information of the bridge (which element, and its global consistency) is lost. If one takes the middle way and
absorbs the discrepancy by blurring, the blurring needed is $\ell$ times the wall and grows with $Q$, while the upper
bound $C$ of IUT~IV does not. On the six completely factored data it is still below $C$ (between $0.36$ and $0.92$ of
$C$); after replacing $q$ by $q^N$ it exceeds $C$ for $N\ge3$ (checked for $N\le40$ prime to $\ell$ with the values of
$C$ in \cite{TC}), and then IUT~IV does not bound it and the comparison says nothing about the height. This is the same
shape as Scholze--Stix's blurring \q{rendering the inequality thus obtained useless} \cite[\S2.2]{SS}.

\subsection{In the language of HoTT}
\tagH{} In the language of \cite{SL}: what can be carried along a path of identification is a property (a
proposition), and the simplification of Scholze--Stix is always correct for propositions, by path induction. If a
structure (a chosen element) is carried along an identification consisting of ``all isomorphisms'', only an orbit of
the automorphism group is determined (Proposition~\ref{p:fullpoly}). An additive bridge is information on the side of
structure, so it cannot become a lax 2-cell that crosses the $\Theta$-link. In function fields the Kodaira--Spencer
morphism could be a bridge because the comparison was made by a connection over a single base $S$, that is, by a
specific isomorphism $p_1^*\Fcal\simeq p_2^*\Fcal$, compatible with the ring structure (Proposition~\ref{p:ks}~(i)).
Mochizuki also describes the incompatibility of the $\Theta$-link with arithmetic degrees as the monodromy of a loop.
\srcquote{\tagS{} \q{Alternatively, one may understand the incompatibility of the $\Theta$-link (i.e., a horizontal
arrow in the log-theta-lattice) with arithmetic degrees (cf. (AD)) as a sort of nontrivial monodromy around the loop
(i.e., a special case of `$\Gamma_{\mathrm{loop}}$'!) that arises as soon as one identifies the domain and codomain
of the $\Theta$-link --- an incompatibility, i.e., `contradiction', that is in fact avoided in IUTch (cf. (IUAD))
precisely by distinguishing the (ring structures in the) domain and codomain of the $\Theta$-link, i.e.,
by working with `$\Gamma_{\mathrm{seg}}$', and then applying the multiradial algorithms of IUTch to obtain an
alternative way to compute this nontrivial monodromy (cf. (GIUT)).}}{Mochizuki, \emph{Report on discussions \dots{}
March 15--20, 2018}, \S14, (LbEx1) \cite{R18}}

\subsection{What HoTT contributed}
\tagH{} Every proposition of this note can be proved without HoTT, and we wrote them in ordinary mathematics.
Proposition~\ref{p:fullpoly} is an elementary lemma of category theory, and Mochizuki's full poly-isomorphisms and
indeterminacies are this idea from the start. The language of HoTT was nevertheless useful, because it puts the
identification, down to ``along which path'', into the types. In HoTT an identification is itself a datum, a path, and
a proposition cannot be written down without naming the path. From an identification that does not choose a path (all
isomorphisms), only properties can be extracted. Applying this distinction, one sees at once that an additive bridge is
information on the side of structure and does not cross the $\Theta$-link. Since the dispute over IUT is about how
mutually alien copies are identified, a type discipline of exactly this kind applies. \cite[\S1.1]{SL} gave three
reasons for choosing HoTT: Mochizuki himself explains the dispute by segment, loop, universal cover and monodromy; an
identification is a path, a first-class object; and univalence decides, as a theorem, when identifying isomorphic
objects is harmless. This note is one example in which those reasons applied.

\tagH{} There are limits as well. We used only the low levels of HoTT (propositions, sets, groupoids), at which the
statements of HoTT translate directly into the language of orbits, torsors and groupoids. The advantage of HoTT is that
it makes the discipline of groupoids the default and enforces it, not that it can say what other languages cannot. The
claims at the level of HoTT are stated on paper and have not been checked by machine (the verification in Lean in
\cite{SL} covers only the quantitative core, and the equality of Lean is incompatible with univalence). This is one
case, and a case favourable to HoTT, since identification itself is what is disputed. One can say that HoTT has
produced new mathematics when a theorem comes out that would have been hard to find without it; for instance, Anel,
Biedermann, Finster and Joyal generalized the Blakers--Massey theorem to an arbitrary higher topos and an arbitrary
modality and write that \q{Our proof is inspired by the one given in Homotopy Type Theory} \cite{ABFJ}. This note
contains no result of that kind.

\section{The dichotomy, and what a repair needs}

\subsection{Table}
\begin{table}[ht]\centering\small
\caption{Bridge and size (the reading of this note).}\label{t:dichotomy}
\begin{tabularx}{\textwidth}{@{}>{\raggedright\arraybackslash}p{3.3cm}>{\raggedright\arraybackslash}X
>{\raggedright\arraybackslash}X>{\raggedright\arraybackslash}X@{}}
\toprule
situation & bridge (that a 2-cell exists) & size (blurring) & shape of the conclusion\\
\midrule
function fields (Kodaira--Spencer) & automatic from the structure ($\kappa$ is a linear morphism) &
$\deg\omega(\log)$, determined by the curve alone & $2\deg\omega_E\le\deg\omega(\log)$\\
curves over number fields with a global multiplicative subspace & the isogeny (an algebraic map) & of the order of
$\log l$ & of the form $l\cdot h\le h+\log l$ (\cite[\S2.3]{Alien}); excluded in IUT by Definition~3.1~(c)
(Proposition~\ref{p:sl2})\\
Baker's method, Stewart--Yu & automatic from the addition (Proposition~\ref{p:additive}) & determined by the lower
bound for linear forms in logarithms, large & a weak abc (a shift in the exponent)\\
Pasten's arithmetic derivatives & exists unconditionally (Lemma~3.5) & $\|\psi\|$; abc $\Rightarrow$ small $\Rightarrow$ abc with some exponent (Corollary~4.6) &
$c/\log c\le\mathrm{rad}(abc)\|\psi\|/\log2$\\
IUT (Step (xi) of Corollary~3.12) & not derived in the written argument (\S6) & no size, since no bridge is derived;
the blurring is bounded, $E\le C$ (IUT~IV) & $(k-1)Q\le E$\\
\bottomrule
\end{tabularx}
\end{table}

\tagH{} In one sentence, the reading of this note is the following. \textbf{IUT gave up the ring structure (in
particular, the addition) in order to circumvent the absence of a global multiplicative subspace, but that ring
structure was precisely the material of a Kodaira--Spencer-type bridge.} The bridge of function fields came from a
specific isomorphism over a single base, compatible with the ring structure (Proposition~\ref{p:ks}~(i), \S7.3). The
introduction of IUT~I names, as the price of simulating a global multiplicative subspace, the \q{dismantling} of the
global prime structure of the number fields involved \cite{IUT1}, and according to \cite[\S3.3, (iv)]{Alien} the operation
of discarding primes can only be performed if one forgets the additive structure (the quotations in \S5). The curves on
which the isogeny bridge acts are excluded by Definition~3.1~(c) (Proposition~\ref{p:sl2}), and an additive bridge does
not cross the $\Theta$-link, which glues multiplicative data by full poly-isomorphisms (Propositions~\ref{p:fullpoly}
and~\ref{p:gap}). What remains is the inner bridge of \S8.2. This is a reading, not a theorem.

\subsection{What a repair needs}
\tagH{} To complete IUT as an argument of Kodaira--Spencer type, one needs the following.
\begin{itemize}
\item \textbf{A bridge built inside the multiradial representation.} The way of bringing in an additive relation from
outside is closed (Propositions~\ref{p:fullpoly} and~\ref{p:gap}). The isogeny bridge is excluded by Definition~3.1~(c)
(Proposition~\ref{p:sl2}).
\item \textbf{That its blurring does not grow with $q$} ($E\le C$), uniformly for all initial $\Theta$-data.
\end{itemize}
\tagH{} For a single datum this claim is equivalent to the height inequality (\cite[Theorem~5.3]{SL}), and if it holds
for all genuine initial $\Theta$-data, the abc conjecture follows (\lean{classicalABC\_of\_forall\_xi} of \cite{lana};
for arithmetic derivatives, the abc conjecture implies the existence of small ones, and their existence implies the
abc conjecture with some exponent \cite[Corollary~4.6]{Pa}). The language of HoTT can state the required 2-cell precisely, but cannot construct it
(\cite[Theorem~5.7]{SL}). That is, the repair amounts not to an adjustment of IUT but to finding a new proof of the
height inequality. One form of it would be to show the (FrgInv) of the source, \q{the operation of forgetting discussed
in (MnAlyLb) can in fact, if one allows for suitable indeterminacies, be inverted} \cite[\S3.11]{EssLgc}, at the level
of values rather than of structures, with an indeterminacy independent of $q$.

\tagS{} Anabelian reconstruction itself constructs a number field group-theoretically from a fundamental group.
Theorem~1.9 of \cite{AbsTopIII} reconstructs, for a curve over a sub-$p$-adic field $k$, the NF-portion of the function
field from the extension $1\to\Delta_X\to\Pi_X\to G_k\to1$. Remark~1.9.4 says that, for a local field $k$ (an MLF), of
the two dimensions of $G_k$ the one represented by the inertia group is not rigid, but that both dimensions are
rigidified once coupled with $\Delta_X$.
\srcquote{\tagS{} \q{there exists a functorial `group-theoretic' algorithm} \dots{} \q{for reconstructing the
`NF-portion of the function field' of $X$ from the extension of profinite groups $1\to\Delta_X\to\Pi_X\to
G_k\to1$} \dots{} \q{On the other hand, the dimension of $G_k$ represented by the inertia group in $I_k\subseteq G_k$
[which, as is well-known, is of cohomological dimension 1] is `far from rigid'} \dots{} \q{The `group-theoretic'
algorithm of Theorem 1.9 shows that the condition of being `coupled with $\Delta_X$' [i.e., via the extension
determined by $\Pi_X$] has the effect of rigidifying both of the 2 dimensions of $G_k$}}{Mochizuki, \emph{Topics in
absolute anabelian geometry III}, Theorem~1.9 and Remark~1.9.4 \cite{AbsTopIII}}
\tagH{} This reconstruction works inside one theater (where $\Pi_X$ is available) and works in the same way for every
curve. lana-agents/iut contains a formalization in Lean of the steps of this reconstruction that recover the additive
structure (Uchida's lemma, Proposition~1.3, and Theorem~1.9~(e); \lean{Iut/Reconstruction/}) \cite{lana}. Across the
$\Theta$-link, only the Galois group $G_v$ and the unit portions (not $\Pi_v$) are carried locally, and only the
realified degrees globally; so the ``coupling'' of Remark~1.9.4 is not available, and the addition returns only up to
the indeterminacy of (Ind2), as we read it. What a repair needs is a \textbf{quantitative version} of this reconstruction
(a claim, measured by log-volumes, of how precisely the addition returns after crossing the $\Theta$-link). It is not
written.

\subsection{Rescue in a weak sense}
\begin{itemize}
\item \tagH{} \textbf{A weak Corollary~3.12 with a large blurring.} If one allows the blurring of the size of Pasten or
Stewart--Yu, a weak form of the comparison holds. But that is only a restatement of existing theorems.
\item \tagH{} \textbf{Curves with a global multiplicative subspace.} Isogenies give a bridge. But, as
Proposition~\ref{p:sl2} shows, these curves lie outside the domain of IUT.
\item \tagH{} \textbf{Function fields.} The Kodaira--Spencer morphism gives a bridge (Proposition~\ref{p:ks}). This is
classical, and Mochizuki himself writes it as the limiting case $N\to\infty$ of the argument (of the type of IUT) that
multiplies the height by $N$ (\cite[\S3.1, (v)]{Alien}).
\end{itemize}
\tagH{} None of these adds new arithmetic. As another kind of rescue, there is the view that Corollary~3.12 is not
``nonsense'' but an exact proposition, a lax 2-cell, whose truth is merely unresolved (\cite{SL}).

\section{Limitations}
\begin{itemize}
\item \textbf{An essay that proposes a viewpoint.} Proposition~\ref{p:ks} is classical, and
Propositions~\ref{p:additive}, \ref{p:legendre}, \ref{p:sl2}, \ref{p:fullpoly} and~\ref{p:gap} are elementary; there is
no new theorem. The dichotomy of Table~\ref{t:dichotomy} is an analogy, not a theorem. The correspondence between
function fields and number fields follows Mochizuki's analogy \cite{Alien} and the analogy of Vojta and Pasten
\cite{Pa}.
\item \textbf{Section~7 is about bridges brought in from outside.} It is a check along the definition of the
$\Theta$-link (full poly-isomorphisms, multiplicative gluing); it does not show that there is no construction that
carries global consistency in another form inside the multiradial representation. If there were one, it would be the
repair of \S8.2.
\item \textbf{Condition (b) of Proposition~\ref{p:gap} is not something IUT requires.} IUT places no additive relation
on the $\Theta$ side. What Proposition~\ref{p:gap} shows is only that a repair bringing in an additive relation from
outside is either contradictory or needs a blurring that grows with $Q$.
\item \textbf{Observation~\ref{o:curve} is a reading on the statements of theorems.} We read the statements of
Theorem~3.11, IUT~II Corollary~4.8, IUT~I Example~5.1 and Remark~3.1.7, but not all of the proofs and the other remarks.
\item \textbf{Literature.} We have not read the papers of Faltings, Stewart--Yu and Kim, \cite{GenEll}, or the text of
Baker's method themselves (we rely on the introduction of Dupuy--Zureick-Brown and on the bibliographies and summaries
of Pasten and \cite{Alien}). Joshi's Construction~V is stated to be in preparation, and we have not confirmed its
existence or contents. The work of Anel--Biedermann--Finster--Joyal is cited from its abstract.
\item \textbf{Data:} four abc triples and seven initial $\Theta$-data; D2357 is conditional. Proposition~\ref{p:legendre}
is a statement about places of multiplicative reduction, and at all places counted $\mathrm{ord}(q)$ was even (checked by
the script).
\item \textbf{The part on HoTT} (\S7.3, \S7.4) depends on the framework and reading of \cite{SL}; the assessment of
\S7.4 is about this one case.
\item No claim is made about whether IUT is correct.
\end{itemize}

\section{Checks and reproduction}
\begin{itemize}
\item \texttt{scripts/bridge\_lib.py} computes Table~\ref{t:triples} (four triples: the valuation of the $p$-adic
logarithm, computed from the series, agrees with the exponent) and Table~\ref{t:data} (seven data: $\mathrm{ord}(q)$ is
even, the gaps, and the label-averaged sum is $\ell$ times the wall). The values of the wall and of $Q$ agree with the
output of the programs of \cite{TC}. Its output is \texttt{outputs/bridge.txt}
(\texttt{python3 scripts/bridge\_lib.py > outputs/bridge.txt}); \texttt{paper/make\_generated.py} writes the tables and
the numbers of the text from the same functions.
\item Data: \texttt{data/D*.txt}, the same files as in \texttt{data/} of \cite{TC} (version 1.0.0).
\item Quotations are taken from the English originals (the text layers of the PDFs, with mathematical symbols and
decorations checked on the page images).
\end{itemize}

\section*{Use of generative AI}
The question of this note --- whether the number-field version of the Kodaira--Spencer argument, which gives the
height inequality for function fields, can be realized as a quantity on the construction of IUT --- and the decision to
write it up as a separate note arose in the author's discussions with generative AI tools; the author judged them
important and decided to make them the subject of this note. Generative AI tools (large language models) were used,
under the author's direction, in the computations, in checking the literature and in the writing. The author takes full
responsibility for the contents of this note.

\begin{thebibliography}{AbsTopIII}\raggedright
\bibitem[ABFJ]{ABFJ}
Mathieu Anel, Georg Biedermann, Eric Finster and Andr\'e Joyal, \emph{A generalized {B}lakers--{M}assey theorem},
J. Topol. \textbf{13} (2020), no.~4, 1521--1553, \href{https://doi.org/10.1112/topo.12163}{doi:10.1112/topo.12163};
arXiv:1703.09050.

\bibitem[AbsTopIII]{AbsTopIII}
Shinichi Mochizuki, \emph{Topics in absolute anabelian geometry {III}: global reconstruction algorithms}, J. Math. Sci.
Univ. Tokyo \textbf{22} (2015), 939--1156.

\bibitem[Alien]{Alien}
Shinichi Mochizuki, \emph{The mathematics of mutually alien copies: from {G}aussian integrals to inter-universal
{T}eichm\"uller theory}, in: Inter-universal Teichm\"uller Theory Summit 2016, RIMS K\^oky\^uroku Bessatsu \textbf{B84}
(2021), 23--192.

\bibitem[DZB]{DZB}
Taylor Dupuy and David Zureick-Brown, \emph{Deligne--{I}llusie classes as arithmetic {K}odaira--{S}pencer classes},
J. Th\'eor. Nombres Bordeaux \textbf{31} (2019), no.~2, 371--383,
\href{https://doi.org/10.5802/jtnb.1086}{doi:10.5802/jtnb.1086}; arXiv:1804.00228.

\bibitem[EssLgc]{EssLgc}
Shinichi Mochizuki, \emph{On the essential logical structure of inter-universal {T}eichm\"uller theory in terms of
logical {AND} ``$\wedge$''/logical {OR} ``$\vee$'' relations: report on the occasion of the publication of the four main
papers on inter-universal {T}eichm\"uller theory}, RIMS Preprint 1968, November 2022; version of March 2024.

\bibitem[Fal]{Fal}
Gerd Faltings, \emph{Does there exist an arithmetic {K}odaira--{S}pencer class?}, in: Algebraic geometry: Hirzebruch 70
(Warsaw, 1998), Contemp. Math. \textbf{241}, Amer. Math. Soc., 1999, 141--146,
\href{https://doi.org/10.1090/conm/241/03633}{doi:10.1090/conm/241/03633}.

\bibitem[Fe]{Fe}
Ivan Fesenko, \emph{Arithmetic deformation theory via arithmetic fundamental groups and nonarchimedean
theta-functions, notes on the work of {S}hinichi {M}ochizuki}, Eur. J. Math. \textbf{1} (2015), no.~3, 405--440,
\href{https://doi.org/10.1007/s40879-015-0066-0}{doi:10.1007/s40879-015-0066-0}.

\bibitem[GalKS]{GalKS}
Shinichi Mochizuki, \emph{The {G}alois-theoretic {K}odaira--{S}pencer morphism of an elliptic curve}, preprint, July
2000.

\bibitem[Ge2]{Ge2}
Lukas Geiger, \emph{The {B}ridge-{C}ontract {C}riterion: a form audit for transport arguments, with an instantiation in
the abc literature} (The abc Gap Series, Part~2), version 0.2, Zenodo, 16 August 2026,
\href{https://doi.org/10.5281/zenodo.21960477}{doi:10.5281/zenodo.21960477}; all versions
\href{https://doi.org/10.5281/zenodo.21960476}{doi:10.5281/zenodo.21960476}.

\bibitem[Ge3]{Ge3}
Lukas Geiger, \emph{The {B}ridge {C}ertificate: a line-level audit instrument for inter-theoretic transfer arguments,
with four applications to the abc literature} (The abc Gap Series, Part~3), version 0.1, Zenodo, 13 August 2026,
\href{https://doi.org/10.5281/zenodo.21925804}{doi:10.5281/zenodo.21925804}; all versions
\href{https://doi.org/10.5281/zenodo.21925803}{doi:10.5281/zenodo.21925803}.

\bibitem[GenEll]{GenEll}
Shinichi Mochizuki, \emph{Arithmetic elliptic curves in general position}, Math. J. Okayama Univ. \textbf{52} (2010),
1--28.

\bibitem[HA]{HA}
Shinichi Mochizuki, \emph{The {H}odge--{A}rakelov theory of elliptic curves: global discretization of local {H}odge
theories}, preprint, September 1999.

\bibitem[Hi]{Hi}
Joaquim Reizi Higuchi, \emph{A formal negative result for the passage from {S}tep (xi-e) to {S}tep (xi-f) in {IUT}
{III}}, preprint, 31 March 2026, \href{https://doi.org/10.13140/RG.2.2.10989.76006}{doi:10.13140/RG.2.2.10989.76006}
(a record of the same day titled `iut bridge preprint' has
\href{https://doi.org/10.13140/RG.2.2.24411.53286}{doi:10.13140/RG.2.2.24411.53286}).

\bibitem[IUT1]{IUT1}
Shinichi Mochizuki, \emph{Inter-universal {T}eichm\"uller theory {I}: construction of {H}odge theaters}, Publ. Res.
Inst. Math. Sci. \textbf{57} (2021), 3--207, \href{https://doi.org/10.4171/PRIMS/57-1-1}{doi:10.4171/PRIMS/57-1-1}.

\bibitem[IUT2]{IUT2}
Shinichi Mochizuki, \emph{Inter-universal {T}eichm\"uller theory {II}: {H}odge--{A}rakelov-theoretic evaluation}, Publ.
Res. Inst. Math. Sci. \textbf{57} (2021), 209--401, \href{https://doi.org/10.4171/PRIMS/57-1-2}{doi:10.4171/PRIMS/57-1-2}.

\bibitem[IUT3]{IUT3}
Shinichi Mochizuki, \emph{Inter-universal {T}eichm\"uller theory {III}: canonical splittings of the log-theta-lattice},
Publ. Res. Inst. Math. Sci. \textbf{57} (2021), 403--626,
\href{https://doi.org/10.4171/PRIMS/57-1-3}{doi:10.4171/PRIMS/57-1-3}.

\bibitem[IUT4]{IUT4}
Shinichi Mochizuki, \emph{Inter-universal {T}eichm\"uller theory {IV}: log-volume computations and set-theoretic
foundations}, Publ. Res. Inst. Math. Sci. \textbf{57} (2021), 627--723,
\href{https://doi.org/10.4171/PRIMS/57-1-4}{doi:10.4171/PRIMS/57-1-4}.

\bibitem[JoQA]{JoQA}
Kirti Joshi, \emph{Questions and answers regarding my preprints on Arithmetic Teichmuller Spaces}, 1 April 2024
(author's web page), Question~2.14.

\bibitem[Kim]{Kim}
Minhyong Kim, \emph{Geometric height inequalities and the {K}odaira--{S}pencer map}, Compositio Math. \textbf{105}
(1997), 43--54, \href{https://doi.org/10.1023/A:1017995416618}{doi:10.1023/A:1017995416618}; Erratum, ibid. \textbf{121}
(2000), 219.

\bibitem[LANA]{LANA}
Project LANA, \emph{Project LANA interim report on IUT theory}, 20 July 2026,
\url{https://ncatlab.org/nlab/files/LANAProject-Report-July2026.pdf}.

\bibitem[lana]{lana}
lana-agents, \emph{iut}: a Lean~4 formalization of inter-universal {T}eichm\"uller theory, GitHub repository,
\url{https://github.com/lana-agents/iut}, commit \texttt{553076a} (3 November 2026). Files cited:
\lean{Iut/\allowbreak{}WallEquiv/\allowbreak{}NonDerivable.lean},
\lean{Iut/\allowbreak{}Reconstruction/\allowbreak{}Uchida/\allowbreak{}FunctionField.lean},
\lean{Iut/\allowbreak{}Reconstruction/\allowbreak{}Thm19/\allowbreak{}NFAdditive.lean},
\lean{Iut/\allowbreak{}Abc/\allowbreak{}FunctionFieldSzpiro.lean}, \lean{Iut/\allowbreak{}Abc/\allowbreak{}P24Table.lean}.

\bibitem[LM]{LM}
R.~Laniewski and K.~M\"uller, \emph{Pasten's arithmetic derivatives and the decomposition into gains},
arXiv:2610.26490v1 [math.NT], 19 October 2026.

\bibitem[Pa]{Pa}
Hector Pasten, \emph{Arithmetic derivatives through geometry of numbers}, Canad. Math. Bull. \textbf{65} (2022), no.~4,
906--923, \href{https://doi.org/10.4153/S0008439521000990}{doi:10.4153/S0008439521000990}; arXiv:2106.16165v3.

\bibitem[Pj]{Pj}
Paul Pajo, \emph{Scale-free common-container comparisons do not entail numerical inequalities: a minimum-viable no-go
theorem motivated by the {IUT} {III} {T}heorem 3.11-to-{C}orollary 3.12 bridge}, preliminary research note,
26 April 2026, \href{https://doi.org/10.13140/RG.2.2.33620.16005}{doi:10.13140/RG.2.2.33620.16005}.

\bibitem[R18]{R18}
Shinichi Mochizuki, \emph{Report on discussions, held during the period March 15--20, 2018, concerning inter-universal
{T}eichm\"uller theory ({IUTch})}, February 2019.

\bibitem[RF]{RF}
Toshihisa Urahigashi, \emph{No reading avoids the height inequality: 36 readings of Corollary 3.12 and the conservation
of the height load}, version 1.0.0, Zenodo, 2026,
\href{https://doi.org/10.5281/zenodo.23965952}{doi:10.5281/zenodo.23965952}.

\bibitem[Roots]{Roots}
Shinichi Mochizuki, \emph{Classical roots of inter-universal {T}eichm\"uller theory}, slides, November 2020, 27 pp.,
\url{https://www.kurims.kyoto-u.ac.jp/~motizuki/2020-11\%20Classical\%20roots\%20of\%20IUT.pdf}.

\bibitem[SL]{SL}
Toshihisa Urahigashi, \emph{Segment or loop? A homotopy-type-theoretic overview of the dispute over Corollary 3.12 of
inter-universal {T}eichm\"uller theory}, version 1.3.0, Zenodo, 2026,
\href{https://doi.org/10.5281/zenodo.23991545}{doi:10.5281/zenodo.23991545}.

\bibitem[SS]{SS}
Peter Scholze and Jakob Stix, \emph{Why abc is still a conjecture}, manuscript, 2018.

\bibitem[StYu]{StYu}
C. L. Stewart and Kunrui Yu, \emph{On the abc conjecture, {II}}, Duke Math. J. \textbf{108} (2001), no.~1, 169--181,
\href{https://doi.org/10.1215/S0012-7094-01-10815-6}{doi:10.1215/S0012-7094-01-10815-6}.

\bibitem[TC]{TC}
Toshihisa Urahigashi, \emph{Two computations of the height inequality and the multiradial algorithm: reading IUT from
the viewpoint of computer science}, version 1.0.0, Zenodo, 2026,
\href{https://doi.org/10.5281/zenodo.23968299}{doi:10.5281/zenodo.23968299}.

\bibitem[TK]{TK}
Takkun-kohinata, \emph{IUT\_LEAN}, Lean repository, \url{https://github.com/Takkun-kohinata/IUT_LEAN}, commit
\texttt{9d56c46} (15 July 2026); file \lean{REQUEST\_theta\_hull\_derivation.md}.

\bibitem[W311]{W311}
Toshihisa Urahigashi, \emph{Where Theorem 3.11 sits: the log-volume content of the results cited in the proof of IUT
III, Corollary 3.12, does not imply its conclusion}, Zenodo, 2026,
\href{https://doi.org/10.5281/zenodo.24028832}{doi:10.5281/zenodo.24028832} (all versions).
\end{thebibliography}

\end{document}
